\chapter{Tools Evaluated and Not Adopted}
\label{ch:rejected}

Each of these was investigated far enough to make a judgement rather than
dismissed on first impression. One was suggested by the client.

\section{GENESYS Excel Connector}
\label{sec:excel}

\subsection{What it does}

The Excel Connector is a GENESYS add-in that moves data between a project and a
spreadsheet. It can download any entity, attribute, parameter or relationship
from a project into Excel, and upload edited values back into the repository.

\subsection{Why it is useful}

Engineers already use Excel, so it meets people where they are rather than asking
them to learn an interface. It supports bulk editing, which is substantially
faster than changing attributes one at a time through the GENESYS interface, and
it keeps the spreadsheet and the repository synchronised rather than producing a
dead export.

The client recommended looking at it, and the recommendation was a good one. It
is the closest existing thing to what this project is building: a bidirectional
connector between GENESYS and an external representation, with a round trip
rather than a one-way export.

\subsection{Why it does not work for this project}

Three reasons, in descending order of weight.

\paragraph{It flattens the structure.}
The connector presents entities and their attributes as rows and columns. What
the translation needs is the shape of the data: which function calls which, which
state transitions to which, how a class relates to its methods. Relationships are
the substance of the model, and a spreadsheet is the wrong representation for a
graph. The tool would be reconstructing structure that the Client API already
hands over directly.

\paragraph{It is not automatable.}
Driving it would mean a script operating the add-in through its user interface,
pressing buttons and waiting for loading, which is fragile and slow. The Client
API performs the same reads and writes programmatically and without the
intermediate file.

\paragraph{It adds a format with no benefit.}
Every translation would pass through a spreadsheet that neither end wants, purely
as a transport. The Client API is available, so there is no reason to accept
that.

\subsection{What it is useful for instead}

It is a good reference for what a connector should do, and the team treats it as
one. It demonstrates the round trip, the synchronisation behaviour and the scope
of what the vendor considers reasonable for an integration to touch. It is a model
to learn from rather than a component to build on.

\section{GENESYS REST API}

Covered in Section~\ref{sec:genesys-api}, since it is best understood alongside
the interface that replaced it. In short: it exists, it would have been the
better choice, and it is not available to the team because enabling it requires
administrative access that the team does not have and the client does not
provide.

\section{Direct Database Access}

\subsection{What it does}

GENESYS stores its repositories in SQL Server, with the database files under the
installation's data directory. The database can be opened with standard tooling
such as SQL Server Management Studio and inspected or modified directly.

\subsection{Why it might be useful}

It bypasses every API. Anything the repository holds is reachable, with no
framework constraint, no licensing component, and no dependence on which
operations the vendor chose to expose. For a project that needs both read and
write access to model content, that is superficially attractive.

\subsection{Why it does not work for this project}

\paragraph{The schema is undocumented and private.}
GENESYS uses a schema-driven entity-relationship-attribute model in which entity
classes are themselves data. The physical tables are a generic store rather than
tables named after the concepts in the model, so using them would mean reverse
engineering a metamodel held as data. That is a project in itself.

\paragraph{Writing through it bypasses everything that keeps the model valid.}
Schema validation, relationship integrity, the view generators that produce
diagrams, and the application's own change tracking would all be circumvented. A
write that looked correct in the database could leave the repository
inconsistent.

\paragraph{It is an implementation detail, not an interface.}
Nothing obliges the vendor to keep the schema stable across versions, so a tool
built on it would break on upgrade. A deliverable the client cannot safely
upgrade is not a deliverable.

\paragraph{Holding the database open prevents GENESYS from starting.}
The database files are served by the GENESYS services while the application runs,
so an external connection contends with a live process. In practice the effect is
worse than contention: with an external tool connected, GENESYS fails to reach
its own repository on startup and falls back to prompting for a licence key, as
though the installation were unlicensed. Recovering from that means closing the
external connection and working back through the licensing wizard, which is
disruptive enough on its own to rule the approach out even if the earlier
objections did not.

\subsection{What it was useful for}

Opening the database read-only was informative during investigation. It confirmed
the entity-attribute storage model and made clear why the schema must be read at
run time rather than assumed. It was reconnaissance, not a route.

\section{GNSX and GNUX Interchange Files}

\subsection{What it does}

GENESYS imports and exports its own interchange format, used for project
backups, baselines and change files. A file can hold a whole repository, a
project, or a subset of one, and the import command loads it back.

\subsection{Why it might be useful}

It needs no API, no authentication and no network. A tool could produce an import
file and have the user load it, which would make the write path trivially simple
and would work regardless of what the Client API does or does not permit. It was
the natural fallback had the Client API turned out to be read-only.

\subsection{Why it does not work for this project}

\paragraph{The University Edition uses a different format.}
The edition available to the team imports and exports GNUX rather than GNSX. The
two are not interchangeable, so files produced during development would not load
into a commercial GENESYS installation, and the client could not use the output.

\paragraph{It is not a live connection.}
Every operation would require the user to export a file, run the tool, and import
the result. That is a worse workflow than the one the Client API makes possible,
and it rules out the push and pull behaviour entirely.

\paragraph{Import semantics are unclear.}
Whether an import merges against existing entities by identifier or creates
duplicates determines whether a round trip is even possible, and this was not
established.
